% 1. RESUMEN
\section{Resumen}
\label{sec:resumen}

La segmentación o pipelining es la técnica microarquitectural predominante para aumentar el throughput de procesadores RISC mediante la superposición temporal de instrucciones. Aunque la literatura recoge numerosas implementaciones de procesadores RISC-V RV32IM en FPGA, tanto monociclo como con pipeline de cinco etapas, los estudios existentes evalúan en su mayoría una única microarquitectura o emplean plataformas y protocolos de medición heterogéneos, de modo que no existe evidencia empírica directa y reproducible que contraste ambas variantes del mismo núcleo bajo condiciones controladas.
Este proyecto aborda esa brecha mediante el diseño e implementación en SystemVerilog de un procesador RISC-V RV32IM en dos microarquitecturas, monociclo y pipeline de cinco etapas con resolución de hazards de datos y control, sobre la FPGA Intel Cyclone V DE1-SoC. La verificación funcional de ambas implementaciones se realiza mediante la suite riscv-tests y testbenches desarrollados con cocotb, cubriendo el conjunto de instrucciones RV32IM completo. A partir de síntesis homogéneas con Intel Quartus Prime y un protocolo de medición unificado, se comparan experimentalmente la frecuencia máxima, el CPI efectivo, el throughput y la utilización de recursos lógicos, memoria y bloques DSP. Adicionalmente, se desarrolla un módulo de visualización de los estados internos del pipeline integrable en la misma tarjeta, concebido como herramienta de apoyo docente para cursos de arquitectura de computadores.
Se espera que la microarquitectura segmentada alcance un throughput superior al del diseño monociclo equivalente gracias a una frecuencia de operación más alta, a pesar del incremento en CPI efectivo introducido por los hazards. Los resultados obtenidos constituirán evidencia cuantitativa y replicable sobre los trade-offs reales entre ambas microarquitecturas en una plataforma FPGA de capacidad media ampliamente utilizada en entornos educativos.

\textbf{Palabras clave:} RISC-V, RV32IM, pipeline, monociclo, FPGA, DE1-SoC, SystemVerilog, hazards, verificación funcional, rendimiento.
